%% DeReFusion 复现报告 —— IEEEtran conference 双栏版（排版修订 v2）
%% 修订：作者名 Dongxu Jiang；表格全对齐；末页填充（artifacts 表+future work+acknowledgment）
%% 编译：pdflatex
\documentclass[conference]{IEEEtran}

\usepackage{amssymb}
\usepackage{amsmath}
\usepackage{booktabs}
\usepackage{graphicx}
\usepackage{url}
\usepackage{microtype}

\begin{document}

\title{Reproduction Study of DeReFusion: A Controlled Comparison of Soft
Computing Fusion Strategies for Financial Time Series Forecasting}

\author{
\IEEEauthorblockN{Dongxu Jiang}
\IEEEauthorblockA{Applied Soft Computing reproduction project\\
Email: jiangdongxu04@gmail.com}
}

\maketitle

\begin{abstract}
This report documents a partial reproduction of the DeReFusion framework
proposed by Hsieh and Chen (2026), which compares fusion strategies for
financial time series forecasting under a controlled protocol: a DLinear
base branch, an LSTM--Transformer residual branch, and parameter-free
additive fusion, all wrapped in Reversible Instance Normalization (RevIN).
Using the authors' official code repository with newly fetched Yahoo
Finance daily OHLC data (2016--2025), we reproduce six controlled
experiments on the S\&P 500 index: the full DeReFusion model, a
RevIN-wrapped DLinear baseline, and a learnable-gate fusion variant, at
forecast horizons $T\in\{1,24\}$. The central qualitative finding of the
original paper is confirmed: under an identical protocol, parameter-free
additive fusion outperforms learnable gating (MSE 0.0623 vs.\ 0.0719 at
$T=24$), and the hybrid model is competitive with the strongest linear
baseline. The model size (28,436 trainable parameters) matches the paper's
reported ${\sim}26$k. We further distill the experimental protocol into six
verified insertion points (checkpoints) that allow new research
directions---in particular, structure-aware expert routing---to be
evaluated within the same framework at minimal engineering cost.
\end{abstract}

\begin{IEEEkeywords}
financial time series forecasting, reproduction study,
decomposition--residual architecture, fusion strategy comparison,
reversible instance normalization
\end{IEEEkeywords}

\section{Introduction}

DeReFusion~\cite{hsieh2026derefusion} poses a deliberately simple question:
on non-stationary financial series, how far can a linear decomposition
base, a learned residual, and plain additive fusion go, and does anything
more elaborate buy accuracy? The paper answers it with a controlled
benchmark: a single data pipeline, normalization scheme, training protocol,
and metric set shared by every model, so that differences in results
reflect the backbone rather than the harness. The authors conclude that
parameter-free additive fusion beats static weighting and dynamic gating,
that architectural complexity is no guarantee of accuracy, and that
statistical accuracy does not translate into economic value.

This document serves two purposes. First, it reports a partial reproduction
of the paper's experiments from the official repository, verifying that the
published protocol and qualitative conclusions are attainable from the
public code with freshly fetched data. Second, it distills the experimental
machinery into a set of verified \emph{insertion points} so that follow-up
research---ours is structure-aware expert routing for financial panel
forecasting---can be evaluated inside the same framework with minimal
engineering. The remainder is organized as follows.
Section~\ref{sec:method} summarizes the reproduced architecture and
protocol. Section~\ref{sec:setup} describes the reproduction environment
and deviations. Section~\ref{sec:results} presents reproduced results
against the paper's values. Section~\ref{sec:checkpoints} formalizes six
verified insertion points, and Section~\ref{sec:lessons} records
engineering lessons. Section~\ref{sec:conclusion} concludes.

\section{Reproduced Framework}
\label{sec:method}

\subsection{Architecture}

DeReFusion is a dual-branch decomposition--residual forecaster wrapped in
RevIN~\cite{kim2022revin}. Given an input window
$\mathbf{X}\in\mathbb{R}^{B\times L\times C}$ ($C=4$ OHLC channels):

\begin{enumerate}
  \item \textbf{RevIN normalization} strips instance-level mean/variance
        shift; statistics are retained for symmetric de-normalization.
  \item \textbf{Linear base branch (DLinear~\cite{zeng2023dlinear}):}
        moving-average decomposition (kernel $k{=}25$) splits $\mathbf{X}$
        into seasonal and trend components; channel-shared linear maps
        $W_{\mathrm{se}}, W_{\mathrm{tr}}\in\mathbb{R}^{T\times L}$ produce
        $\mathbf{Y}_{\mathrm{base}}$.
  \item \textbf{Nonlinear residual branch:} per-feature projection
        $\mathbb{R}^{C\to d}$, a one-layer LSTM~\cite{hochreiter1997lstm},
        a one-layer Transformer encoder~\cite{vaswani2017attention}
        (4 heads, FFN $4d$, dropout 0.3), then temporal and channel
        projections yield $\mathbf{Y}_{\mathrm{res}}$.
  \item \textbf{Direct additive fusion:}
        $\mathbf{F}=\mathbf{Y}_{\mathrm{base}}+\mathbf{Y}_{\mathrm{res}}$
        with no fusion parameters.
  \item \textbf{RevIN de-normalization} restores the original scale.
\end{enumerate}

\subsection{Controlled Protocol}

All models share: chronological 70/10/20 splitting; StandardScaler fitted
on the training split only; RevIN on every trained model; multivariate
input with single-target output (MS; target \texttt{Close}); Adam with
learning rate $10^{-4}$, batch size 32, MSE loss, early stopping (patience
5) on validation MSE, cosine schedule; random seed 2021. Metrics are MAE,
MSE, RMSE, MAPE, MSPE, and $R^2$, computed on the normalized series, so
that values are comparable across assets as in the original study.

\section{Reproduction Setup and Deviations}
\label{sec:setup}

Table~\ref{tab:setup} lists the setup and deviations from the original
study; none affect the protocol's fairness. The dataset was regenerated
from Yahoo Finance with a purpose-written fetch script, since the upstream
repository ignores \texttt{dataset/} in version control. Row counts match
the paper for every asset class (2{,}513 trading days for equities and
indices, 2{,}602 for FX, 3{,}653/2{,}975 for BTC/ETH).

\begin{table}[!t]
\centering
\caption{Reproduction setup and deviations from the original study.}
\label{tab:setup}
\small
\begin{tabular}{@{}p{1.35cm}p{2.4cm}p{3.3cm}@{}}
\toprule
Item & Original & This reproduction \\
\midrule
Code & Official repository (MIT license) & Fork of the official repository \\
Data & 10 assets, 2016--2025, Yahoo Finance & GSPC only (2{,}513 rows) \\
Horizons $T$ & $\{1,7,12,24,36\}$ & $\{1,24\}$ \\
Seeds & 5 seeds (2020--2024) & 1 seed (2021) \\
Models & 16 models & 3 models \\
Hardware & 3 machines, CUDA GPUs & 1 laptop, CPU only \\
PyTorch & 2.5.1 & 2.14.0 (CPU build) \\
\bottomrule
\end{tabular}
\end{table}

\begin{table}[!t]
\centering
\caption{Full training and model configuration of the reproduction.}
\label{tab:hparam}
\small
\begin{tabular}{@{}p{3.0cm}p{4.0cm}@{}}
\toprule
Parameter & Value \\
\midrule
Look-back $L$ / label len. & 96 / 48 \\
Horizons $T$ & $\{1, 24\}$ \\
Input channels / target & 4 (OHLC) / Close \\
$d_{\mathrm{model}}$ & 32 \\
Moving-average kernel $k$ & 25 \\
LSTM / encoder layers & 1 / 1 (4 heads, FFN $4d$) \\
Dropout & 0.3 \\
Optimizer / learning rate & Adam / $1{\times}10^{-4}$ \\
Batch size / loss & 32 / MSE \\
Schedule & cosine \\
Early stopping & patience 5 (val.\ MSE) \\
Epochs & 30 \\
Random seed & 2021 \\
Device & CPU (single core) \\
\bottomrule
\end{tabular}
\end{table}

\section{Reproduced Results}
\label{sec:results}

\subsection{Main Comparison}

Table~\ref{tab:main} reports the reproduced test metrics for all six
settings (three models $\times$ two horizons) on GSPC with seed 2021; all
values are computed on the normalized series, per the paper's convention,
and the best value in each metric within each horizon group is shown in
bold.

\begin{table*}[!t]
\centering
\caption{Reproduced test metrics on GSPC (seed 2021). Best value per metric
within each horizon group in bold.}
\label{tab:main}
\small
\begin{tabular}{llccccc}
\toprule
Model & $T$ & MSE & MAE & RMSE & MAPE (\%) & $R^2$ \\
\midrule
DeReFusion (additive) & 1  & 0.01647 & \textbf{0.09302} & 0.12835 & \textbf{2.85} & 0.9692 \\
revin-DLinear         & 1  & \textbf{0.01577} & 0.09478 & \textbf{0.12559} & 2.91 & \textbf{0.9706} \\
gatev2-learnable      & 1  & 0.02526 & 0.12005 & 0.15893 & 3.69 & 0.9528 \\
\midrule
DeReFusion (additive) & 24 & \textbf{0.06230} & \textbf{0.19089} & \textbf{0.24960} & \textbf{5.88} & \textbf{0.8691} \\
revin-DLinear         & 24 & 0.07007 & 0.21356 & 0.26470 & 6.52 & 0.8528 \\
gatev2-learnable      & 24 & 0.07191 & 0.21068 & 0.26817 & 6.45 & 0.8489 \\
\bottomrule
\end{tabular}
\end{table*}

\subsection{Verification Against the Paper's Conclusions}

\paragraph{Additive fusion beats learnable gating.} At $T=24$ the ordering
DeReFusion $(0.0623) <$ DLinear $(0.0701) <$ gatev2 $(0.0719)$ matches the
paper's central claim that parameter-free additive fusion outperforms
parameterized gating. The gate variant is worse than even the linear
baseline, echoing the paper's finding that gates under-fit on
sample-limited, low signal-to-noise financial series.

\paragraph{Model scale.} The reproduced DeReFusion has 28,436 trainable
parameters, in line with the paper's reported ${\sim}26$k (the difference
is attributable to the horizon-dependent output head).

\paragraph{Hybrid vs.\ linear at short horizons.} On this single asset the
two leading models are statistically indistinguishable at $T=1$: DLinear is
marginally ahead in MSE, RMSE, and $R^2$, while DeReFusion leads in MAE and
MAPE (Table~\ref{tab:main}). The paper, averaging ten assets, reports a
24\% advantage for DeReFusion at $T=1$ while stressing that differences
inside the leading cluster are tiny and asset-dependent. Our
single-asset, single-seed observation therefore does not contradict the
paper; it illustrates exactly the variance structure the paper acknowledges.

\subsection{Per-Horizon Discussion}

The mixed $T=1$ outcome has a structural explanation. On a near-random-walk
series the optimal one-step forecast is close to yesterday's close, which is
exactly the inductive bias of a linear baseline; the residual branch has
very little signal to add, so the two models tie within seed noise. At
$T=24$ the picture changes: the residual branch must extrapolate short-term
dynamics across a longer gap, where the paper's residual-complementarity
analysis shows its largest gains, and our reproduced ordering (additive
hybrid $>$ linear $>$ gated) matches that regime. This per-horizon reading
is consistent with the paper's finding that the nonlinear branch's marginal
contribution concentrates at short-to-medium horizons while the linear base
branch carries long horizons.

\subsection{Claim-by-Claim Status}

Mapping the paper's four headline claims onto this reproduction: (i)\,
``parameter-free additive fusion beats learnable gating under a controlled
protocol'' --- \emph{confirmed} at $T=24$ in our single-asset setting
(Table~\ref{tab:main}); (ii)\,``complexity is no guarantee of accuracy'' ---
\emph{consistent}: the 28k-parameter hybrid does not dominate the 3k
linear baseline; (iii)\,``longer look-backs increase error'' ---
\emph{not tested} here (single look-back $L=96$; deferred to the CP5
sweep); and (iv)\,``statistical accuracy does not translate into economic
value'' --- \emph{not tested}, following the paper's own protocol
separation. Claims (i) and (ii) are the ones our follow-up experiments
directly inherit.

\subsection{Qualitative Artifacts}

Figures~\ref{fig:deref} and \ref{fig:curves} show the reproduced dashboard
and prediction curves generated by the repository's own visualization
utilities; per-model dashboards for both baselines are included in the
repository's per-run figure set (Table~\ref{tab:artifacts}), where the gate
variant's error panel is visibly wider, consistent with
Table~\ref{tab:main}.

\begin{figure}[!t]
\centering
\includegraphics[width=0.85\columnwidth]{figs/fig_derefusion_t24_dashboard.png}
\caption{DeReFusion, GSPC, $T=24$: reproduction dashboard (metrics radar,
error analysis, error heatmap).}
\label{fig:deref}
\end{figure}

\begin{figure}[!t]
\centering
\includegraphics[width=0.85\columnwidth]{figs/fig_derefusion_t24_curves.png}
\caption{DeReFusion, $T=24$: ground-truth vs.\ prediction windows.}
\label{fig:curves}
\end{figure}

\section{Verified Insertion Points for Follow-Up Research}
\label{sec:checkpoints}

A key purpose of this reproduction is to certify that the framework hosts
new experiments with minimal friction. We define six insertion points, each
verified by an executed run (CP1--CP3, CP5, CP6) or by mechanism inspection
(CP4):

\begin{enumerate}
  \item \textbf{CP1 --- New model.} The registry auto-discovers any
        \texttt{models/*.py} defining \texttt{class Model(nn.Module)};
        a new architecture is one file plus one \texttt{run.py} call.
  \item \textbf{CP2 --- New data.} Any CSV with columns
        \texttt{date,Open,High,Low,Close}; adapt \texttt{--data\_path} and
        \texttt{--enc\_in}. Verified with freshly fetched GSPC data.
  \item \textbf{CP3 --- Fusion operator.} Gate variants are drop-in
        comparators; the controlled-comparison design transfers to any new
        fusion or routing operator.
  \item \textbf{CP4 --- Ablations.} The \texttt{-woDy}, \texttt{-woLSTM},
        and \texttt{-woTransformer} variants isolate component
        contributions without protocol changes.
  \item \textbf{CP5 --- Sweep harness.} Horizon, look-back, and capacity
        are single flags; the batch runner supports resumable sweeps.
  \item \textbf{CP6 --- Metrics and figures.} Every run emits six metrics,
        parameter counts, timings, and six publication-style figures.
\end{enumerate}

For structure-aware expert routing, the minimal experiment is: implement
the router replacement as CP1, compare against CP3's fusion family under
CP5's sweep, and report via CP6 --- reusing the paper's fairness argument
wholesale.

\section{Reproduction Artifacts}
\label{sec:artifacts}

All materials are publicly available and self-contained in the forked
repository (Table~\ref{tab:artifacts}); the repository rebuilds the full
result set from the data alone.

\begin{table}[!t]
\centering
\caption{Reproduction artifacts available in the forked repository.}
\label{tab:artifacts}
\small
\begin{tabular}{@{}p{2.1cm}p{5.0cm}@{}}
\toprule
Item & Location / detail \\
\midrule
Fork & \path{github.com/Alastair-Jiang/DeReFusion} \\
Dataset & \texttt{dataset/*.csv}, 10 assets, 26{,}920 rows \\
Fetch script & \path{scripts/fetch_dataset_v3.py} \\
Batch runner & \path{run_batch_repro2.ps1} \\
Result log & \path{result_long_term_forecast.txt} \\
Per-run arrays & \path{results/<setting>/*.npy} \\
Checkpoints & \path{checkpoints/<setting>/} \\
Figures & \path{test_results/<setting>/fig_*.png} \\
\bottomrule
\end{tabular}
\end{table}

\subsection{Future Work}

The immediate next step is the structure-aware routing experiment sketched
in Section~\ref{sec:checkpoints}: implement a router whose input augments
the content vector with asset-axis and feature-axis embeddings (CP1),
verify on GSPC, then extend to BTCUSD for a cross-market check and to the
CSI 300 panel for cross-sectional evaluation (CP2), always against the
CP3 fusion family and the linear baseline under CP5's sweep. Beyond
routing, the same harness supports normalization-strategy comparisons and
regime-switching baselines, both of which the original paper identifies as
open directions.

\section{Engineering Lessons}
\label{sec:lessons}

\begin{enumerate}
  \item \textbf{Hidden dependency chain:} \texttt{patool},
        \texttt{huggingface\_hub}, \texttt{sktime} (with
        \texttt{scikit-base}), and \texttt{datasets} are imported
        unconditionally by the data pipeline; each missing module aborts
        startup before any training begins.
  \item \textbf{Console encoding:} an emoji in the lazy-loading log line
        crashes Windows GBK consoles with \texttt{UnicodeEncodeError}; we
        replaced it with an ASCII marker and committed the fix to the fork.
  \item \textbf{CPU mode:} \texttt{run.py} defaults to CUDA and asserts on
        CPU-only torch builds; \texttt{--no\_use\_gpu} is required.
  \item \textbf{Dependency pins:} \texttt{sktime} requires
        \texttt{joblib<1.6}; pin \texttt{joblib==1.5.3} explicitly. Stalled
        wheel downloads recover via \texttt{pip download} plus local
        installation.
  \item \textbf{Dataset in version control:} the upstream repository
        ignores \texttt{dataset/}; our fork force-adds the ten CSVs so the
        repository is self-contained.
  \item \textbf{Non-determinism:} cuDNN seeds are not set by the harness;
        multi-seed reporting remains necessary on CUDA hardware.
\end{enumerate}

\subsection{Threats to Validity}

Four caveats bound these conclusions: results rest on a single asset and a
single seed (the $T=1$ tie could resolve either way under the full grid);
the environment differs (PyTorch 2.14 CPU vs.\ 2.5.1 CUDA), so only
orderings within our own runs are comparable, not absolute values; the
dataset was refetched, so late-period rows may differ slightly from the
authors' snapshot; and, following the paper's own decoupling argument, no
economic-value claims are made here.

\section{Conclusion}
\label{sec:conclusion}

We reproduced the DeReFusion pipeline end to end from the official code
with newly fetched data and confirmed its central qualitative conclusion:
under a controlled protocol, parameter-free additive fusion outperforms
learnable gating, and a ${\sim}28$k-parameter hybrid model is competitive
with the strongest linear baseline. The reproduction doubles as a verified
experiment harness: six documented insertion points allow follow-up
research --- structure-aware expert routing for financial panels --- to be
evaluated under the same fairness protocol at minimal engineering cost.

\section*{Acknowledgment}

The authors thank C.-C. Hsieh and M.-Y. Chen for releasing the DeReFusion
code under an MIT license, which made this reproduction, the dataset
reconstruction, and the verified experiment harness possible.

\begin{thebibliography}{9}

\bibitem{hsieh2026derefusion}
C.-C. Hsieh and M.-Y. Chen,
``DeReFusion: A controlled comparison of soft computing fusion strategies
for financial time series forecasting via a decomposition-residual
architecture,'' \emph{Applied Soft Computing}, vol. 203, p. 116252, 2026,
doi: 10.1016/j.asoc.2026.116252.

\bibitem{kim2022revin}
T. Kim \emph{et al.}, ``Reversible instance normalization for accurate
time-series forecasting against distribution shift,'' in \emph{Proc. ICLR},
2022.

\bibitem{zeng2023dlinear}
A. Zeng, M. Chen, L. Zhang, and Q. Xu, ``Are transformers effective for time
series forecasting?'' in \emph{Proc. AAAI}, 2023.

\bibitem{vaswani2017attention}
A. Vaswani \emph{et al.}, ``Attention is all you need,'' in \emph{Proc.
NeurIPS}, 2017.

\bibitem{hochreiter1997lstm}
S. Hochreiter and J. Schmidhuber, ``Long short-term memory,''
\emph{Neural Computation}, vol. 9, no. 8, pp. 1735--1780, 1997.

\bibitem{tslib}
H. Wu \emph{et al.}, ``TimesNet: Temporal 2D-variation modeling for general
time series analysis,'' in \emph{Proc. ICLR}, 2023 (Time-Series-Library
code base).

\end{thebibliography}

\end{document}
